\section{Results}
\label{sec:results}
No independently verified result exists. The blinded review that the protocol requires has not been performed, and the main evaluation of Section~\ref{sec:design} has not been run. This section reports what was obtained instead, each part labelled by its basis. Subsections~\ref{sec:res-samples} to~\ref{sec:res-planning} are \emph{exploratory and benchmark-derived}: they come from the offline replay of two small real samples, and ``agreement with the benchmark'' means agreement with the provisional ORCID label of Section~\ref{sec:benchmark}, not identity accuracy. No figure here is pooled with another, corrected by an audit, or called precision of the system. Subsection~\ref{sec:res-pending} lists the protocol results that are not reported.

\subsection{Samples and Scope}
\label{sec:res-samples}
\textbf{Smoke sample.} 50 cases from 23 blocks, captured with 92 hashed responses (Section~\ref{sec:smoke}). It is not a probability sample, and its cases are related within blocks. We therefore give every figure for the 50 cases and, for the selective results, also for one case per block (23 cases).

\textbf{Exploratory sample.} A second sample of 50 blocks, one primary case each, was drawn from the main frame with the derived main seed of freeze F4 (strata S, M and L with 23, 24 and 3 cases), disjoint from the smoke blocks and the 250 pilot blocks. No credential was available and no money was to be spent, so the capture used OpenAlex's free keyless tier, with the stop guard of Section~\ref{sec:main} set to a cost ceiling of USD 0.10. After 38 successful requests the provider answered HTTP 429, because the shared free daily budget was exhausted, and the capture stopped after the 18th case (nine S and nine M cases from 18 different blocks). The cases are processed in the order of their opaque identifiers, which is unrelated to the strata, so these 18 are a random subset of the drawn 50; they are reported as an \emph{18-case exploratory sample} with no inference to the frame. An earlier attempt on this sample sent 70 requests that were all recorded as decoding errors because of a defect in the stop guard (it re-wrapped an already decoded body with the wire response's compression header, which mock transports without compression could not show); that run produced no usable data, was kept as a negative result, and the defect was fixed with a regression test. Intervals below are exact Clopper--Pearson intervals that treat cases as independent; for the smoke sample they are optimistic because of clustering.

\subsection{Outcome Categories}
\label{sec:res-outcomes}
% generated by scripts/analyze_free_evaluation.py (T019); do not edit by hand
\begin{table*}[t]
\caption{Outcome categories of the offline replay under each condition and policy (benchmark-independent counts). Smoke: 50 cases from 23 blocks. Exploratory: 18 complete cases from 18 blocks. Auto = automatic, Man = manual review, Unres = unresolved, Empty = successful empty pool, Fail = provider failure, NoSc = no defined score.}
\label{tab:exp-outcomes}
\centering
\footnotesize
\begin{tabular}{@{}llrrrrrrrrrrrr@{}}
\toprule
 & & \multicolumn{6}{c}{Smoke (50)} & \multicolumn{6}{c}{Exploratory (18)} \\
\cmidrule(lr){3-8}\cmidrule(lr){9-14}
Cond. & Policy & Auto & Man & Unres & Empty & Fail & NoSc & Auto & Man & Unres & Empty & Fail & NoSc \\
\midrule
C0 & selective & 0 & 50 & 0 & 0 & 0 & 0 & 0 & 18 & 0 & 0 & 0 & 0 \\
C0 & gates+P3 & 0 & 50 & 0 & 0 & 0 & 0 & 0 & 18 & 0 & 0 & 0 & 0 \\
C0 & forced top-1 & 50 & 0 & 0 & 0 & 0 & 0 & 18 & 0 & 0 & 0 & 0 & 0 \\
C1 & selective & 4 & 46 & 0 & 0 & 0 & 0 & 0 & 18 & 0 & 0 & 0 & 0 \\
C1 & gates+P3 & 4 & 46 & 0 & 0 & 0 & 0 & 0 & 18 & 0 & 0 & 0 & 0 \\
C1 & forced top-1 & 50 & 0 & 0 & 0 & 0 & 0 & 18 & 0 & 0 & 0 & 0 & 0 \\
C1x & selective & 7 & 43 & 0 & 0 & 0 & 0 & 3 & 15 & 0 & 0 & 0 & 0 \\
C1x & gates+P3 & 26 & 24 & 0 & 0 & 0 & 0 & 8 & 10 & 0 & 0 & 0 & 0 \\
C1x & forced top-1 & 50 & 0 & 0 & 0 & 0 & 0 & 18 & 0 & 0 & 0 & 0 & 0 \\
\bottomrule
\end{tabular}
\end{table*}

Table~\ref{tab:exp-outcomes} counts the six outcome categories under every condition and policy. In every configuration and both samples each case was either automatic or manual: there was no unresolved, empty-pool, provider-failure or no-score outcome. Under C0 (name only) no case was assigned automatically by either policy with safeguards or by the threshold gates, as the policy guarantees by construction. Under the primary condition C1 with the full selective policy, 4 of the 50 smoke cases (8\%, 95\% CI 2--19\%) and 0 of the 18 exploratory cases (upper limit 19\%) were assigned automatically. The threshold-gate ablation gave the same counts under C1 (4 and 0), so the safeguards B1 to B5 changed no outcome under C1 in these samples. Under the exploratory condition C1x, whose publication evidence comes from a member of the identity group and is therefore leakage-prone, the full policy assigned 7 of 50 and 3 of 18 cases, while the ablation without the safeguards assigned 26 of 50 and 8 of 18: the safeguards withheld automation in 19 and 5 cases respectively.

\subsection{Candidate Retrieval}
\label{sec:res-retrieval}
% generated by scripts/analyze_free_evaluation.py (T019); do not edit by hand
\begin{table*}[t]
\caption{Benchmark-derived retrieval diagnostics (provisional ORCID agreement, not identity accuracy). The pool is the captured R2 pool; R0 is descriptive and shares its response with R2 whenever the requests are identical. Rank is the provider's order. Intervals are Clopper--Pearson and ignore clustering of cases within blocks.}
\label{tab:exp-retrieval}
\centering
\footnotesize
\begin{tabular}{@{}lrrrr@{}}
\toprule
 & \multicolumn{2}{c}{Smoke (50)} & \multicolumn{2}{c}{Explor.\ (18)} \\
\cmidrule(lr){2-3}\cmidrule(lr){4-5}
Measure (R2 pool) & Count & 95\% CI & Count & 95\% CI \\
\midrule
Expected ORCID in pool & 38/50 & [0.62, 0.87] & 13/18 & [0.47, 0.90] \\
Success@1 (end to end) & 22/50 & [0.30, 0.59] & 7/18 & [0.17, 0.64] \\
Success@5 & 34/50 & [0.53, 0.80] & 12/18 & [0.41, 0.87] \\
Success@10 & 38/50 & [0.62, 0.87] & 13/18 & [0.47, 0.90] \\
Conditional Recall@1 & 22/38 & [0.41, 0.74] & 7/13 & [0.25, 0.81] \\
R0: expected ORCID in pool & 36/50 & [0.58, 0.84] & 12/18 & [0.41, 0.87] \\
\bottomrule
\end{tabular}
\end{table*}

Table~\ref{tab:exp-retrieval} gives the retrieval diagnostics. The R2 pool never failed and was never empty, with a mean of 6.1 candidates in the smoke sample and 8.6 in the exploratory sample (at most ten in both). The profile carrying the benchmark ORCID was in the pool for 38 of 50 smoke cases (76\%, CI 62--87\%) and for 13 of 18 exploratory cases (72\%, CI 47--90\%); it was the provider's first candidate in 22 of 50 (44\%) and in 7 of 18 (39\%), and within the first five in 34 of 50 and 12 of 18. Among the cases whose pool contained it, the provider ranked it first in 22 of 38 and 7 of 13. The end-to-end mean reciprocal rank was 0.54 and 0.50. R0 and R2 differed little: R0 found the profile in 36 of 50 and 12 of 18 cases, and the only difference was that R0 returned an empty pool in 2 and 1 cases, which R2's initials fallback filled. Since the two share their primary response whenever the requests are identical, no comparison of the two strategies is made. For 12 of 50 and 5 of 18 cases the expected ORCID was absent from a successful pool. These counts include cases in which the right person's profile carries no ORCID, which the diagnostic cannot see, so they bound label availability and not retrieval accuracy.

\subsection{Selective Resolution (Benchmark-Derived)}
\label{sec:res-selective}
% generated by scripts/analyze_free_evaluation.py (T019); do not edit by hand
\begin{table*}[t]
\caption{Benchmark-derived selective results. Coverage is the share of cases assigned automatically; precision is the share of automatic assignments whose selected profile carries the benchmark ORCID (provisional, never independent verification). Forced top-1 never abstains, so its precision is a top-1 agreement rate. With zero errors in $n$ automatic assignments the one-sided 95\% upper bound on the error rate is $1-0.05^{1/n}$.}
\label{tab:exp-selective}
\centering
\footnotesize
\begin{tabular}{@{}llrrrr@{}}
\toprule
Dataset & Cond. / policy & Automatic & Coverage (95\% CI) & Agree with benchmark & Precision (95\% CI) \\
\midrule
Smoke, 50 cases & C1, selective & 4 & 8\% [0.02, 0.19] & 3 & 75\% [0.19, 0.99] \\
Smoke, 50 cases & C1, gates+P3 & 4 & 8\% [0.02, 0.19] & 3 & 75\% [0.19, 0.99] \\
Smoke, 50 cases & C1, forced top-1 & 50 & 100\% [0.93, 1.00] & 17 & 34\% [0.21, 0.49] \\
Smoke, 50 cases & C1x, selective & 7 & 14\% [0.06, 0.27] & 7 & 100\% [0.59, 1.00] \\
Smoke, 50 cases & C1x, gates+P3 & 26 & 52\% [0.37, 0.66] & 25 & 96\% [0.80, 1.00] \\
\midrule
Smoke, 23 blocks & C1, selective & 2 & 9\% [0.01, 0.28] & 2 & 100\% [0.16, 1.00] \\
Smoke, 23 blocks & C1, gates+P3 & 2 & 9\% [0.01, 0.28] & 2 & 100\% [0.16, 1.00] \\
Smoke, 23 blocks & C1, forced top-1 & 23 & 100\% [0.85, 1.00] & 8 & 35\% [0.16, 0.57] \\
Smoke, 23 blocks & C1x, selective & 4 & 17\% [0.05, 0.39] & 4 & 100\% [0.40, 1.00] \\
Smoke, 23 blocks & C1x, gates+P3 & 12 & 52\% [0.31, 0.73] & 11 & 92\% [0.62, 1.00] \\
\midrule
Exploratory, 18 & C1, selective & 0 & 0\% [0.00, 0.19] & 0 & undefined \\
Exploratory, 18 & C1, gates+P3 & 0 & 0\% [0.00, 0.19] & 0 & undefined \\
Exploratory, 18 & C1, forced top-1 & 18 & 100\% [0.81, 1.00] & 6 & 33\% [0.13, 0.59] \\
Exploratory, 18 & C1x, selective & 3 & 17\% [0.04, 0.41] & 3 & 100\% [0.29, 1.00] \\
Exploratory, 18 & C1x, gates+P3 & 8 & 44\% [0.22, 0.69] & 7 & 88\% [0.47, 1.00] \\
\bottomrule
\end{tabular}
\end{table*}

Table~\ref{tab:exp-selective} gives coverage and benchmark agreement. Under C1 with the full policy, 3 of the 4 automatic smoke assignments carried the benchmark ORCID (75\%, exact CI 19--99\%), and 2 of 2 in the one-case-per-block view; the exploratory sample had no automatic assignment, so no precision is defined. The one disagreeing smoke assignment cannot be classified without verification: it may be a wrong assignment or a benchmark labelling problem. Under C1x the automatic assignments agreed with the benchmark in 7 of 7 and 3 of 3 cases under the full policy, but C1x is exploratory and leakage-prone. The forced top-one policy, which never abstains and uses HAMMER's score ranking, agreed with the benchmark for 17 of 50 smoke cases (34\%) and 6 of 18 exploratory cases (33\%) under C1, whereas the provider's own first candidate agreed for 22 of 50 and 7 of 18. HAMMER's heuristic ordering of the pool therefore did not improve on the provider's order in these samples.

\subsection{Why Coverage Is Low: Gate Failures}
\label{sec:res-gates}
Of the 68 cases, all had a candidate pool under C1, and the top candidate's score was high: the median was 0.88 in both samples and only one case had a score below the automatic threshold of 0.75. The binding conditions were competition and evidence. Among the 50 smoke cases, 37 had a margin below 0.15 against a scored competitor, 44 a dominance ratio below 1.35, and 14 fewer than two observed signals for the top candidate; among the 18 exploratory cases the counts were 14, 17 and 3 (the conditions overlap). The institution signal was observed for the top candidate in 29 of 50 and 11 of 18 cases, and candidates of one pool were scored over different signal sets in 36 of 50 and 10 of 18 cases, the situation that safeguard B4 treats. In short, a high score is common, but it is rarely separated from the scores of similarly named competitors.

\subsection{Planning the Verified Evaluation}
\label{sec:res-planning}
% generated by scripts/analyze_free_evaluation.py (T019); do not edit by hand
\begin{table}[t]
\caption{Main-sample size $N$ needed for at least $m'$ automatic assignments with probability 0.95, by the stratified predictive rule, under \emph{assumed} pilot outcomes (coverage $c$ in every stratum, pilot of 100, 100 and 50 cases; frame 7,938, 8,260, 1,123 and 9 blocks, XL a census). This is a planning illustration and not data.}
\label{tab:exp-planning}
\centering
\footnotesize
\begin{tabular}{@{}rrrrr@{}}
\toprule
$c$ & $m'$ & $N$ (Jeffreys) & $N$ (uniform) & $N$ (lower bound) \\
\midrule
0.05 & 16 & 465 & 411 & 1478 \\
0.05 & 32 & 951 & 850 & $>$2500 \\
0.05 & 100 & $>$2500 & $>$2500 & $>$2500 \\
0.10 & 16 & 221 & 206 & 487 \\
0.10 & 32 & 435 & 412 & 876 \\
0.10 & 100 & 1359 & 1294 & 2439 \\
0.20 & 16 & 110 & 106 & 190 \\
0.20 & 32 & 209 & 203 & 339 \\
0.20 & 100 & 630 & 618 & 937 \\
\bottomrule
\end{tabular}
\end{table}

The label-free automatic rate under C1 was 4 of 68 cases (6\%) across both samples, which are not a probability sample, so we use it only to pick a planning range. Table~\ref{tab:exp-planning} shows the main-sample size that the stratified predictive rule would need under assumed pilot outcomes. For a target of $m'=32$ automatic assignments, about 435 primary cases are needed at a coverage of 0.10 and about 950 at 0.05, and a target of 100 assignments is out of reach within $N_{\max}=2{,}500$ at 0.05. A verified evaluation at the observed coverage would therefore have to review every one of its roughly 30 automatic assignments and also audit a sample of several hundred cases, which is a reviewer workload the project has not yet arranged.

\subsection{Protocol Results Not Reported}
\label{sec:res-pending}
The protocol of Section~\ref{sec:design} yields four results that need the label-free coverage pilot, the main sample and independent reviewers: a case-accounting table for the frozen main evaluation, the independently verified precision of automatic assignments, the audit of benchmark labels, and reviewer agreement. None of these has been run, so none is reported here, and the benchmark-derived observations above must not be read as substitutes for them. Independent verification, confirmatory sampling and reviewer-agreement analysis are left to future work, and no independently verified automatic precision or reviewer-agreement estimate is claimed in this study.
